\chapter{Pruebas, Implementación y Resultados}

\section{Implementación y reproducción del entorno}

La infraestructura declarada en Docker Compose está compuesta por tres servicios conectados mediante una red bridge. Sus contenedores se nombran \path{archivum_app}, \path{archivum_db} y \path{archivum_nginx}. El procedimiento para construir e iniciar los servicios es:

\begin{verbatim}
docker compose up -d --build
\end{verbatim}

La configuración monta el código de la aplicación, mantiene los datos de MariaDB en un volumen y comparte los recursos compilados de Vite. La aplicación espera el estado saludable de la base de datos antes de ejecutar migraciones. Para reproducir el entorno en Windows se requiere Docker Desktop; en Linux, Docker Engine y el complemento Docker Compose. Estas instrucciones describen la configuración disponible; no se reporta una validación independiente en cada sistema operativo si no se ejecutó allí.

\section{Control de versiones y catálogo compartido}

El código fuente se mantiene en un repositorio Git alojado en GitHub. Para compartir el catálogo y sus imágenes entre instalaciones, el comando \texttt{catalogo:exportar-seeder} prepara los datos del seeder curado; después, el equipo puede incorporar esos recursos versionados y ejecutar:

\begin{verbatim}
php artisan db:seed --class=CatalogoCuradoSeeder
\end{verbatim}

La prueba de sincronización comprueba que la exportación genera los archivos esperados, que el seeder restaura registros e imágenes y que volver a ejecutarlo no duplica determinados datos. La forma de colaboración del equipo no se describe mediante convenciones de ramas o mensajes de commit que no estén verificadas en el repositorio.

\section{Pruebas funcionales y automatizadas}

El repositorio incluye pruebas PHPUnit para servicios, controladores, acceso por roles, autenticación, catálogo, inventario, compras, exportaciones, devoluciones y sincronización de seeders. La siguiente tabla resume la cobertura funcional identificada en las suites; el resultado de ejecución se actualizará con la salida real del comando de pruebas.

\begin{table}[htbp]
\centering
\caption{Áreas cubiertas por las pruebas automatizadas existentes.}
\label{tab:pruebas_modulos}
\small
\begin{tabularx}{\textwidth}{|p{3.1cm}|X|}
\hline
\textbf{Área} & \textbf{Comprobaciones representativas} \\ \hline
Catálogo y checkout & Categorías, productos, carrito, pedido y seguimiento simulado. \\ \hline
Inventario & Resolución de sucursal, existencias, movimientos y reglas de stock. \\ \hline
Compras & Alta de productos, validaciones de datos e imágenes. \\ \hline
Semillas de catálogo & Exportación, restauración de imágenes e idempotencia. \\ \hline
Autenticación y acceso & Recuperación de contraseña, autenticación y permisos de roles. \\ \hline
Devoluciones y reportes & Procesamiento de devoluciones y exportación de reportes. \\ \hline
\end{tabularx}
\normalsize
\end{table}

La ejecución global se realiza con:

\begin{verbatim}
php artisan test --compact
\end{verbatim}

Las pruebas permiten verificar reglas de negocio y respuestas funcionales en el entorno de pruebas. No se consideran pruebas de carga ni una demostración de concurrencia real. Tampoco se informan tiempos de respuesta, porcentajes de mejora o tasas de reducción de errores sin mediciones reproducibles.

\section{Resultados}

El resultado verificable del desarrollo es una aplicación con módulos enlazados a rutas, controladores, modelos y vistas, además de migraciones y pruebas para los procesos señalados. El catálogo curado y sus recursos gráficos se pueden exportar y restaurar mediante herramientas propias del proyecto. El checkout conserva un flujo simulado con reservas temporales; las operaciones de inventario registran movimientos asociados a existencias y sucursales.

La ejecución de \texttt{php artisan test --compact} concluyó con 190 pruebas aprobadas, cero fallidas y 1\,335 afirmaciones. Este resultado corresponde a las pruebas automatizadas del repositorio en su entorno de pruebas; no representa una medición de rendimiento, una validación de despliegue en múltiples plataformas ni una demostración concurrente de ausencia de sobreventa.

\section{Conclusiones}

Archivum integra en una sola aplicación web un catálogo especializado y funciones administrativas para inventario, proveedores, ventas, pedidos, devoluciones y seguimiento. El diseño relacional separa los atributos generales del producto de las especificaciones de libros, discos y coleccionables, y relaciona variantes SKU con existencias y movimientos.

La infraestructura declarada con Docker Compose facilita reproducir los servicios de aplicación, base de datos y servidor web. Laravel organiza las rutas, modelos, controladores, middleware, vistas y migraciones; las pruebas PHPUnit cubren módulos funcionales relevantes. El alcance actual debe entenderse con dos precisiones: el checkout es simulado y las recepciones a proveedores se almacenan como movimientos, no como órdenes de compra persistentes en una tabla propia.

\section{Trabajos futuros}

\begin{enumerate}
    \item Incorporar una integración de pago externa con confirmación y conciliación de transacciones.
    \item Diseñar un módulo de órdenes de compra con partidas, estados, autorización y vínculo explícito con recepciones.
    \item Añadir integración con servicios de paquetería para generar guías y consultar estados.
    \item Ampliar la validación con pruebas concurrentes para reservas y ajustes de existencias.
    \item Medir rendimiento y usabilidad con escenarios reproducibles y criterios cuantitativos previamente definidos.
    \item Evaluar el soporte de facturación fiscal electrónica y las aplicaciones móviles para operaciones de almacén.
\end{enumerate}
